\section{模型评价与结论}
\subsection{模型特点}
本文以同一计算图表示为基础，按三个场景逐步扩展状态与约束：场景A把每个子图视为独立Task并令跨图数据经DDR；场景B引入同核复用和跨核同步；场景C进一步显式记录共享L2的FIFO状态与带宽竞争。这样的分层建模保留了问题之间可比较的划分、映射和顺序决策，又避免将某场景不存在的通信或缓存收益误用于另一场景。

求解采用结构化初解、关键路径优先调度与有界局部修复。连通分量、共享输入、双计算管道负载和缓存压力共同决定候选的结构；关键链等待与资源空隙决定局部调度；官方评估程序负责最终合法性和完成时间裁决。通过保留原方案、冻结预算及记录逐组结果，算法可区分预测收益与实际收益，适合在100图、1至5核的固定配置上复现。

\subsection{局限与改进方向}
该方法是启发式算法，没有已证明的最优下界。局部成本模型虽考虑资源阶段和事件传播，仍可能错误排序候选；对于共享DDR的并发传输、同步链重构和L2替换，局部预测尤其敏感。有限官方评价预算保护了运行时间，却也限制了复杂图的探索范围。问题一的搜索阶段最长约364秒，但大图的冷启动单核基准与最终复核未包含在该数值中，故目前不能宣称所有用例的端到端求解都满足5至10分钟目标。

后续改进应以冻结候选池为前提比较代理排序与官方排序，再针对误差最大的资源事件校准；对关键输出附近采用保持无环的联合边界、映射、顺序调整；并为不同瓶颈结构分配搜索预算。任何新增机制须与原算法同预算、同图、同核数对照，且逐组核验官方完成时间和额外搬运量。若目标转向其他NPU或内存层级，需替换相应执行事件和资源容量，而不能直接沿用本文的数值参数。

\subsection{结论}
在指定100图与固定硬件配置下，本文完成三种场景的500组方案求解与官方评估。场景A的五核平均加速比为3.665903，场景B为4.159974，场景C所选L2方案为4.240375。它们对应不同的通信、同步与缓存规则，不应把三者差值解释成单一算法模块的纯收益。三问共同表明：有效的多核方案取决于切图通信、关键路径等待、核内计算与共享存储竞争的联合平衡；最终判断依据是官方模拟的总完成时间，额外搬运量作为同时间下的次级目标。
